% Part: first-order-logic
% Chapter: introduction
% Section: syntax

\documentclass[../../../include/open-logic-section]{subfiles}

\begin{document}

\olfileid{fol}{int}{syn}

\olsection{Vormregels (syntaxis)}

Eerst moeten we precies afspreken welke reeksen symbolen tellen als
!!{sentence}s van de eerste-ordelogica. De precieze regels komen later;
voorlopig kijken we naar voorbeelden. De woordenschat van de
eerste-ordelogica heeft twee delen. Met het eerste deel wijzen we
bepaalde dingen aan en zeggen we er iets over. Een !!{constant} wijst
een ding aan; met een !!{predicate} zeggen we iets over het ding dat
we hebben aangewezen. Zo kunnen we $\Obj a$ als !!a{constant} gebruiken
om één ding aan te wijzen. Met de !!{sentence}~$\Atom{\Obj P}{\Obj a}$
zeggen we vervolgens iets over dat ding. Als je bij ``$\Obj a$'' en
``$\Obj P$'' al een betekenis in gedachten hebt, kun je
$\Atom{\Obj P}{\Obj a}$ als een Nederlandse zin lezen. Waarschijnlijk
deed je dat al toen je voor het eerst formele logica leerde. Van zulke
eenvoudige !!{sentence}s van de eerste-ordelogica kunnen we
ingewikkeldere zinnen maken. Daarvoor gebruiken we het tweede deel van
de woordenschat: de logische symbolen, namelijk connectieven en
kwantoren. We kunnen bijvoorbeeld uitdrukkingen vormen als
$(\Atom{\Obj P}{\Obj a} \land \Atom{\Obj Q}{\Obj b})$
of~$\lexists[\Obj x][\Atom{\Obj P}{\Obj x}]$.

Voor degelijk metalogisch onderzoek hebben we precieze definities
nodig: van de semantiek en van de regels in onze systemen voor
!!{derivation}. Daarom moeten we eerst exact afspreken wat telt als
!!a{sentence} van de eerste-ordelogica. Het basisidee is eenvoudig.
Sommige eenvoudige !!{sentence}s, zoals~$\Atom{\Obj P}{\Obj a}$, maken
we alleen met !!{predicate}s en !!{constant}s. Met connectieven en
kwantoren maken we daar ingewikkeldere zinnen van. Maar volgens welke
regels mogen we precies ingewikkeldere !!{sentence}s maken? Zonder die
regels hebben we ``!!{sentence} van de eerste-ordelogica'' nog niet
precies genoeg gedefinieerd. We moeten drie dingen regelen. Ten eerste:
welke reeksen symbolen tellen als !!{sentence}s? Ten tweede: hoe leggen
we de regels zo vast dat we wiskundige bewijzen kunnen geven over
\emph{alle} !!{sentence}s? Ten slotte moeten we eenvoudige bewerkingen
op !!{sentence}s precies definiëren. Een voorbeeld is: ``vervang elke
$\Obj x$ in~$!A$ door~$\Obj b$.'' Door de kwantoren en !!{variable}s in
de eerste-ordelogica is niet meteen duidelijk hoe dat moet. Telt
bijvoorbeeld $\lexists[\Obj x][\Atom{\Obj P}{\Obj a}]$ als
!!a{sentence}? En
$\lexists[\Obj x][\lexists[\Obj x][\Atom{\Obj P}{\Obj x}]]$? Wat
krijgen we als resultaat van ``vervang $\Obj x$ door~$\Obj b$ in
$(\Atom{\Obj P}{\Obj x} \land \lexists[\Obj x][\Atom{\Obj P}{\Obj
x}])$''?

\end{document}
